

\subsection{Computing Describing Functions Numerically}
\label{sect:computing-df}
If it is not practical to derive an analytic expression for a
describing function, it could be evaluated numerically by using
a fast Fourier transform and taking the fundamental harmonic.
Repeating this at several oscillation amplitudes (and possibly
frequency), then interpolating the result.
% An experimental framework for doing this requires the user to supply
% 3 functions either as an input file named-section\index{named-section}
% \Sectref{sect:input-file} or in a file. These functions are compiled
% automatically at the beginning of a \Fina job, and run to create a set
% of interpolation coefficients giving the describing function as a
% function of specified generalized coordinate amplitudes. The user then
% applies the describing function to a matrix by calling a \Fina function,
% \Code{dfval}, either in a matrix-element equation or a custom evaluation
% function.  More details are in \Sectref{sect:df-approx}.

% \subsection{Creating a describing function}\label{sect:df-approx}
\index{describing function!creating}
% Analytical forms of many describing functions have been 
% derived \cite{gelb1968multiple}; if an appropriate one is not
% found, it may be possible to compute a numerical approximation
% using \Eqn{eqn:df-defn}.
An experimental framework for doing this requires the user to supply
3 functions either as an input file named-section \Sectref{sect:input-file} or
in a file. These functions are compiled automatically at the beginning
of a \Fina job, and run to create a set of interpolation coefficients
giving the describing function as a function of specified generalized
coordinate amplitudes. The user then applies the describing function
to a matrix by calling a \Fina function, \Code{dfval}, either in
a matrix-element equation or a custom evaluation function. When called
from an equation it has the form
\begin{lstlisting}
dfval(dfid, gcno, {parameter-values})
\end{lstlisting}
so for example if the input file has a section named \Code{mybilinear}
containing the 3 functions needed to define a describing function which
needs 2 parameter values, to use it to multiply the $10^{th}$ diagonal
and give the 2 parameters the values $0.01$ and $3.5$:
\begin{lstlisting}
[10,10] *= dfval(mybilinear, 10, {0.01,3.5})
\end{lstlisting}

As with other builtin functions, when called from within a custom
function an additional parameter is required, and the declaration looks
like
\begin{lstlisting}
Real dfval(pset& ps, const string& dfid, int gcno,
               const vector<double>& params);
\end{lstlisting}

The 3 functions needed to create a describing function
must have the following signatures:
% \vspace{0.5em}
% \verb!vector<double> df_qs(const vector<double>& params)!
% \vspace{-.5em}
\begin{lstlisting}
   vector<double> df_qs(const vector<double>& params)
\end{lstlisting}
\begin{quote}
returns a vector with the oscillation amplitudes desired in the describing function
\end{quote}
% \verb!double df_fq(double y, const vector<double>& params)!
% \vspace{-.5em}
\begin{lstlisting}
   double df_fq(double y, const vector<double>& params)
\end{lstlisting}
\begin{quote}
returns the value of the ``force'' at displacement \Code{y},
\end{quote}
\begin{lstlisting}
   Interp* df_interp(const string& iid,
         const vector<double>& params,
         const vector<double>& qs, const vector<double>& fs)
\end{lstlisting}
% \vspace{-1.5em}
\begin{quote}
returns a new Interp object containing the interpolants used to fit the
describing function values \Code{fs}
\end{quote}
where \Code{params} are parameter values given as an argument to \Code{dfval}.

Examples of this experimental technique for computing describing
functions are in \Code{nl4.fp} \Sectref{ex:nl1234} and \Code{latent.fp}
\Sectref{ex:latent}.
